% !PRINTED XIII
\weberpage{XIII}

\inhaltabschnitt{Sechzehnter Abschnitt.}{Klassenzahl in quadratischen Körpern.}
\inhaltpar{110}{Fundamentale Einheiten in den Ordnungen}{398}
\inhaltpar{111}{Die \gesperrt{Dirichlet}sche Grenzformel}{402}
\inhaltpar{112}{Klassenzahl}{405}
\inhaltpar{113}{Die Anzahl der Geschlechter}{409}

\inhaltbuch{Drittes Buch.}{Komplexe Multiplikation.}

\inhaltabschnitt{Siebzehnter Abschnitt.}{Elliptische Funktionen und quadratische Formen.}
\inhaltpar{114}{Singuläre Perioden der doppelt periodischen Funktionen}{413}
\inhaltpar{115}{Die singulären Werte der Invariante $j(\omega)$}{418}
\inhaltpar{116}{Klassenzahlrelationen}{423}
\inhaltpar{117}{Arithmetische Natur der Klassenfunktion $H_m(u)$}{426}
\inhaltpar{118}{Komposition der quadratischen Formen}{428}
\inhaltpar{119}{Die Diskriminante der Invariantengleichung}{431}

\inhaltabschnitt{Achtzehnter Abschnitt.}{\gesperrt{Galois}sche Gruppe der Klassengleichung.}
\inhaltpar{120}{Relationen zwischen den Klasseninvarianten derselben Diskriminante}{435}
\inhaltpar{121}{Trennung der entgegengesetzten Klassen}{437}
\inhaltpar{122}{Irreducibilität}{442}
\inhaltpar{123}{Beziehungen zwischen den Klasseninvarianten in den verschiedenen Ordnungen}{450}
\inhaltpar{124}{Klassenkörper und Ordnungskörper}{455}

\inhaltabschnitt{Neunzehnter Abschnitt.}{Berechnung der Klasseninvarianten.}
\inhaltpar{125}{Die Klasseninvariante $\gamma_2$}{457}
\inhaltpar{126}{Die Klasseninvarianten $f(\omega)^{24}$}{462}
\inhaltpar{127}{Die Potenzen von $f(\omega)$ als Klasseninvarianten}{467}
\inhaltpar{128}{Die ersten Fälle der Berechnung von $f(\sqrt{-m})$}{474}
\inhaltpar{129}{Anwendung der Transformation zweiter Ordnung zur Berechnung von Klasseninvarianten}{476}
\inhaltpar{130}{Berechnung von Klasseninvarianten aus den \gesperrt{Schlaefli}schen Modulargleichungen}{477}
\inhaltpar{131}{Berechnung von Klasseninvarianten aus den irrationalen Formen der Modulargleichungen}{485}
